
\documentclass{article}
\usepackage{colortbl}
\usepackage{makecell}
\usepackage{multirow}
\usepackage{supertabular}

\begin{document}

\newcounter{utterance}

\centering \large Interaction Transcript for game `bbh', experiment `formal\_fallacies', episode 226 with reflexion{-}Qwen3.5{-}27B.
\vspace{24pt}

{ \footnotesize  \setcounter{utterance}{1}
\setlength{\tabcolsep}{0pt}
\begin{supertabular}{c@{$\;$}|p{.15\linewidth}@{}p{.15\linewidth}p{.15\linewidth}p{.15\linewidth}p{.15\linewidth}p{.15\linewidth}}
    \# & \multicolumn{2}{c}{Player} && \multicolumn{2}{c}{Game Master} \\
    \hline

    \theutterance \stepcounter{utterance}  
    & & & \multicolumn{4}{p{0.6\linewidth}}{
        \cellcolor[rgb]{0.9,0.9,0.9}{
            \makecell[{{p{\linewidth}}}]{
                \texttt{\tiny{[P1$\langle$GM]}}
                \texttt{You have attempted this same task before and did not succeed. These are the plans you wrote after those attempts.} \\
\texttt{Trial \#1: New plan: I correctly formalized the premises and deduced that the specific individual \$a\$ (who is a sister of Ruth but not a schoolmate of Karen) must be a classmate of Marisa. However, I failed to recognize that this deduction (\$R(a) \textbackslash{}land C(a)\$) directly contradicts the conclusion (\$\textbackslash{}exists x (R(x) \textbackslash{}land \textbackslash{}neg C(x))\$). Since the premises force the existence of a sister of Ruth who *is* a classmate of Marisa, and do not guarantee the existence of one who is *not*, the argument is invalid. My mistake was stopping the analysis before explicitly comparing the derived fact (\$C(a)\$) with the required conclusion (\$\textbackslash{}neg C(x)\$). Next time, after deriving the properties of the existential instance from the premises, I will explicitly check if those properties satisfy the conclusion. If the premises force a condition that contradicts the conclusion (or fail to support it), I will immediately classify the argument as "invalid" and provide the final answer.} \\
\\ 
\\ 
\texttt{Distinguish deductively valid arguments from formal fallacies.} \\
\\ 
\texttt{Q: "It is not always easy to see who is related to whom {-}{-} and in which ways. The following argument pertains to this question: To begin with, Lesley is a close friend of Fernando. Moreover, being a close friend of Fernando or a schoolmate of Lowell is sufficient for being a great{-}grandfather of Leroy. It follows that Lesley is a great{-}grandfather of Leroy."} \\
\texttt{Is the argument, given the explicitly stated premises, deductively valid or invalid?} \\
\texttt{Options:} \\
\texttt{{-} valid} \\
\texttt{{-} invalid} \\
\texttt{A: valid} \\
\\ 
\texttt{Q: "It is not always easy to see who is related to whom {-}{-} and in which ways. The following argument pertains to this question: Whoever is not a great{-}grandfather of Clyde is a stepbrother of Brian. Being an ancestor of Dana is sufficient for not being a great{-}grandfather of Clyde. We may conclude: Everyone who is an ancestor of Dana is a stepbrother of Brian, too."} \\
\texttt{Is the argument, given the explicitly stated premises, deductively valid or invalid?} \\
\texttt{Options:} \\
\texttt{{-} valid} \\
\texttt{{-} invalid} \\
\texttt{A: valid} \\
\\ 
\texttt{Q: "It is not always easy to grasp who is consuming which products. The following argument pertains to this question: Every infrequent user of Paul Mitchell shampoo is either a rare consumer of Nioxin shampoo or a loyal buyer of Caress soap, or both. No regular consumer of Lush soap is a rare consumer of Nioxin shampoo and, in the same time, a loyal buyer of Caress soap. It follows that whoever is an infrequent user of Paul Mitchell shampoo is not a regular consumer of Lush soap."} \\
\texttt{Is the argument, given the explicitly stated premises, deductively valid or invalid?} \\
\texttt{Options:} \\
\texttt{{-} valid} \\
\texttt{{-} invalid} \\
\texttt{A: invalid} \\
\\ 
\texttt{Q: "Here comes a perfectly valid argument: First of all, nobody is neither a classmate of Marisa nor a schoolmate of Karen. Next, not every sister of Ruth is a schoolmate of Karen. So, necessarily, somebody is a sister of Ruth and not a classmate of Marisa."} \\
\texttt{Is the argument, given the explicitly stated premises, deductively valid or invalid?} \\
\texttt{Options:} \\
\texttt{{-} valid } \\
\texttt{{-} invalid} \\
\texttt{A:} \\
            }
        }
    }
     \\ \\

    \theutterance \stepcounter{utterance}  
    & \multicolumn{4}{p{0.6\linewidth}}{
        \cellcolor[rgb]{0.561,0.737,0.561}{
            \makecell[{{p{\linewidth}}}]{
                \texttt{\tiny{[P1$\rangle$GM]}}
                \texttt{To determine whether the argument is valid or invalid, let's break down the premises and the conclusion using formal logic.} \\
\\ 
\texttt{**1. Define the Predicates:**} \\
\texttt{Let the domain be all people.} \\
\texttt{*   \$C(x)\$: \$x\$ is a classmate of Marisa.} \\
\texttt{*   \$S(x)\$: \$x\$ is a schoolmate of Karen.} \\
\texttt{*   \$R(x)\$: \$x\$ is a sister of Ruth.} \\
\\ 
\texttt{**2. Formalize the Premises:**} \\
\\ 
\texttt{*   **Premise 1:** "Nobody is neither a classmate of Marisa nor a schoolmate of Karen."} \\
\texttt{    *   "Neither \$C\$ nor \$S\$" translates to \$\textbackslash{}neg C(x) \textbackslash{}land \textbackslash{}neg S(x)\$.} \\
\texttt{    *   "Nobody is..." means there does not exist an \$x\$ such that...} \\
\texttt{    *   \$\textbackslash{}neg \textbackslash{}exists x (\textbackslash{}neg C(x) \textbackslash{}land \textbackslash{}neg S(x))\$} \\
\texttt{    *   By De Morgan's laws and quantifier negation, this is equivalent to: \$\textbackslash{}forall x (\textbackslash{}neg (\textbackslash{}neg C(x) \textbackslash{}land \textbackslash{}neg S(x))) \textbackslash{}iff \textbackslash{}forall x (C(x) \textbackslash{}lor S(x))\$.} \\
\texttt{    *   **Meaning:** Everyone is either a classmate of Marisa or a schoolmate of Karen (or both).} \\
\\ 
\texttt{*   **Premise 2:** "Not every sister of Ruth is a schoolmate of Karen."} \\
\texttt{    *   "Every sister of Ruth is a schoolmate of Karen" would be \$\textbackslash{}forall x (R(x) \textbackslash{}to S(x))\$.} \\
\texttt{    *   "Not every..." negates this: \$\textbackslash{}neg \textbackslash{}forall x (R(x) \textbackslash{}to S(x))\$.} \\
\texttt{    *   This is equivalent to: \$\textbackslash{}exists x \textbackslash{}neg (R(x) \textbackslash{}to S(x))\$.} \\
\texttt{    *   Since \$A \textbackslash{}to B\$ is equivalent to \$\textbackslash{}neg A \textbackslash{}lor B\$, the negation is \$A \textbackslash{}land \textbackslash{}neg B\$.} \\
\texttt{    *   **Meaning:** There exists at least one person \$a\$ such that \$R(a)\$ is true and \$S(a)\$ is false.} \\
\texttt{    *   **Formal:** \$\textbackslash{}exists x (R(x) \textbackslash{}land \textbackslash{}neg S(x))\$.} \\
\\ 
\texttt{**3. Formalize the Conclusion:**} \\
\\ 
\texttt{*   **Conclusion:** "Somebody is a sister of Ruth and not a classmate of Marisa."} \\
\texttt{    *   **Formal:** \$\textbackslash{}exists x (R(x) \textbackslash{}land \textbackslash{}neg C(x))\$.} \\
\\ 
\texttt{**4. Logical Deduction:**} \\
\\ 
\texttt{1.  From Premise 2, we know there exists a specific individual, let's call them \$a\$, such that:} \\
\texttt{    *   \$R(a)\$ is true (\$a\$ is a sister of Ruth).} \\
\texttt{    *   \$S(a)\$ is false (\$a\$ is NOT a schoolmate} \\
            }
        }
    }
    & & \\ \\

    \theutterance \stepcounter{utterance}  
    & & & \multicolumn{2}{p{0.3\linewidth}}{
        \cellcolor[rgb]{0.9,0.9,0.9}{
            \makecell[{{p{\linewidth}}}]{
                \texttt{\tiny{[GM$|$GM]}}
                \texttt{Target: invalid} \\
            }
        }
    }
    & & \\ \\

    \theutterance \stepcounter{utterance}  
    & & & \multicolumn{2}{p{0.3\linewidth}}{
        \cellcolor[rgb]{0.9,0.9,0.9}{
            \makecell[{{p{\linewidth}}}]{
                \texttt{\tiny{[GM$|$GM]}}
                \texttt{game\_result = LOSE} \\
            }
        }
    }
    & & \\ \\

    \theutterance \stepcounter{utterance}  
    & & \multicolumn{4}{p{0.6\linewidth}}{
        \cellcolor[rgb]{0.95,0.95,0.95}{
            \makecell[{{p{\linewidth}}}]{
                \texttt{\tiny{[P1$\leftrightarrow$P]}}
                \texttt{Post{-}episode Reflexion (not sent to GM):} \\
\texttt{New plan: I will complete the logical deduction by applying Premise 1 to the specific individual \$a\$ identified in Premise 2. Since Premise 2 establishes that \$a\$ is a sister of Ruth (\$R(a)\$) and not a schoolmate of Karen (\$\textbackslash{}neg S(a)\$), and Premise 1 states that everyone is either a classmate of Marisa or a schoolmate of Karen (\$C(x) \textbackslash{}lor S(x)\$), it logically follows that \$a\$ must be a classmate of Marisa (\$C(a)\$). I will then compare this derived fact (\$R(a) \textbackslash{}land C(a)\$) with the argument's conclusion, which claims there exists someone who is a sister of Ruth and *not* a classmate of Marisa (\$\textbackslash{}exists x (R(x) \textbackslash{}land \textbackslash{}neg C(x))\$). Since the premises force the specific sister \$a\$ to be a classmate of Marisa, they do not support the existence of a sister who is *not* a classmate. Therefore, the argument is invalid. I will explicitly state this contradiction and output "invalid" as the final answer.} \\
            }
        }
    }
    & \\ \\

\end{supertabular}
}

\end{document}
